Reinforcement Learning~(RL)-based GraphRAG frameworks, such as Graph-R1~\citep{luo2025graph}, aim to learn an output reasoning trajectory $\boldsymbol{y}$ that leads to a final answer for an input question $\boldsymbol{x}$ by interleaving language reasoning with retrieval over the structured external knowledge graph $\mathcal{G}_K$.

\paragraph{Context similarity-based graph retrieval.}
GraphRAG leverages knowledge graphs~$\mathcal{G}_K$ to model the relational information inherent in the external knowledge.
In this work, we use a hypergraph as external knowledge representation since it expresses $n$-ary relations among more than two entities while conventional graphs model a connection between two entities.
Formally, a hypergraph is defined as $\mathcal{G}_K = \left(V, E\right)$, where $v \in V$ denotes entities and $E$ denotes a set of edges.
Generally, given the knowledge corpus ${K} = \left\{\boldsymbol{d}_1,\boldsymbol{d}_2 \dots, \boldsymbol{d}_N \right\}$, a (hyper)graph is extracted from knowledge documents $\boldsymbol{d}$ through an LLM:
\begin{equation}
    \mathcal{G}_K = \left(V, E \right), \text{where } \left\{\left(e_i, \mathcal{V}_{e_i} \right) \right\}_{i=1}^m \sim \pi_{\text{ext}}\left(\boldsymbol{d}\right)
\end{equation}
where $V, E$ are sets of nodes and hyperedges, respectively.
$\pi_{\text{ext}}$ is an LLM that is prompted to transform the knowledge segment $\boldsymbol{d}$ into a set of relation-entity pairs $\left\{\left(e_i, \mathcal{V}_{e_i} \right) \right\}$, where $e_i$ and $\mathcal{V}_{e_i}=\left\{v_1, \dots, v_n \right\}$ denote the hyperedge and its participating entities, respectively.
On the constructed knowledge graph, conventional methods~\citep{jin2025search} perform the retrieval based on their semantic similarity to the query representation.
It can be formulated as:
\begin{equation}
\begin{split}
&\mathcal{R}_{H}\left(\boldsymbol{q}_t \right) = \operatorname*{argmax}_{e_i \in E}^{k} \Big( \text{sim}(\phi(\boldsymbol{q}_t),\phi(e_i)) \Big), \quad \mathcal{F}_{t}^{\text{ret}} = \bigcup_{e_j \in \mathcal{R}_{H}} \{(e_j, V_{e_j}) \mid  e_j \in E\},
\end{split}
\end{equation}
where $\phi(\boldsymbol{q}_t)$ denotes the query embeddings, and $\phi(e_i)$ represents the hyperedge embedding.
The top-$k$ retrieved set $\mathcal{F}_t^{\text{ret}}$ forms the final knowledge set.
Graph-R1 applies a hybrid hyperedge retrieval strategy where both entity and hyperedge similarity measurements are leveraged to retrieve hyperedges.

\noindent\textbf{Group-Relative Policy Optimization.}
Group-Relative Policy Optimization~(GRPO)~\citep{shao2024deepseekmath, deepseekai2025deepseekv32} is one of the representative RL approaches with its strong performance and efficiency.
We apply GRPO to train the agent, which is formulated as:
\begin{equation}
% \small
\begin{split}
&\mathcal{J}_{\mathrm{GRPO}}(\theta)
= \mathbb{E}_{\left[\boldsymbol{x} \sim \mathcal{D}_X, \{\boldsymbol{y}^{(i)}\}_{i=1}^N \sim \pi_{\theta_{\text{old}}}(\cdot | \boldsymbol{x}; \mathcal{G}_K)\right]}
\\
&\left[
\frac{1}{\left\lvert \boldsymbol{y}^{(i)} \right\rvert}\sum_{t=1}^{\left\lvert \boldsymbol{y}^{(i)} \right\rvert}
\min\left(
\hat{r}_{t}^{(i)}\hat{A}^{(i)},\operatorname{clip}\!\left(\hat{r}_{t}^{(i)},\, 1 -  \epsilon,\, 1 +\epsilon\right)\,\hat{A}^{(i)}
\right)
\right]
-\beta \mathbb{D}_{\text{KL}}\left(\pi_\theta \mid\mid \pi_{\text{ref}} \right),
\end{split}
\end{equation}
where $\hat{r}_{t}^{(i)} = \frac{\pi_\theta\left( {y}^{(i)}_t \mid \boldsymbol{x}, \boldsymbol{y}_{<t}^{(i)}; \mathcal{G}_K\right)}{\pi_\text{old}\left(y_t^{(i)}\mid \boldsymbol{x}, \boldsymbol{y}_{<t}^{(i)}; \mathcal{G}_K\right)}$ denotes the likelihood ratio between current $\pi_\theta$ and the old policy model $\pi_{\text{old}}$. 
The advantage is calculated as $\hat{{A}}^{(i)} =  \frac{R\left(\boldsymbol{x}, \boldsymbol{y}^{(i)}\right)-\text{mean}\left(\left\{R\left(\boldsymbol{x},\boldsymbol{y}^{(j)} \right) \right\}_{j=1}^N \right)}{F\left(\left\{R\left(\boldsymbol{x},\boldsymbol{y}^{(j)} \right) \right\}_{j=1}^N \right)},$
where $F\left(\cdot \right)$ is the normalizer within the group~$\left\{R\left(\boldsymbol{x},\boldsymbol{y}^{(i)}
\right)\right\}_{i=1}^N$.
Most existing approaches apply outcome-supervised settings~\citep{guo2025deepseek}, which assign the same reward at every token in each output $\boldsymbol{y}$ based on the sequence-level reward of the output $\boldsymbol{y}$.

Despite the improvements in reasoning performance of RL-based agentic GraphRAG frameworks, they are still in an early stage and face two key challenges.
(i) the retrieval is largely based on contextual similarity, which may overlook the richer relational and topological structure inherent in the hypergraph.
(ii) They mainly rely on outcome-level supervision, assigning the same trajectory-level reward across all reasoning steps and therefore providing limited credit assignment for the intermediate retrieval and reasoning that are actually beneficial for solving the problem.